% =============================================================================
% Chapitre 3 — Automatisation métier : fonctions, triggers et procédures stockées
% =============================================================================
\chapter{Automatisation métier : triggers et procédures stockées}

\epigraphe{Ce qui doit être fait à chaque fois ne doit dépendre de la mémoire de personne.}

Jusqu'ici, la base \emph{stocke} et \emph{répond}. Dans ce chapitre, elle \emph{agit} : elle refuse
qu'on modifie un contrat, recalcule le statut d'un prêt dès qu'un paiement arrive, inscrit chaque
changement dans un journal, prévient le recouvrement d'un passage en défaut, encaisse les paiements et
calcule les provisions de fin de mois. C'est la partie « automatisation » de ton projet, celle qui fait
la différence entre un exercice de requêtes et un vrai système.

\begin{observation}[Tes fichiers pour ce chapitre]
\begin{itemize}
\item Ton code : \code{sql/04_chapitre3.sql}. Le corrigé complet : \code{corriges/chapitre3.sql}.
\item Prérequis : chapitres 1 et 2 exécutés (tables \code{journal_statuts} et
\code{provisions_mensuelles}, vue \code{v_situation_prets}).
\end{itemize}
\end{observation}

\begin{situation}{Mise en situation --- Le week-end où personne n'est là}
Awa Koné n'a payé aucune mensualité depuis le 25 octobre. Le dimanche 24 janvier 2027, son retard
atteint 91 jours : elle est en défaut au sens réglementaire. Personne n'est au bureau.

Lundi matin, ta responsable veut que le prêt soit déjà \textbf{passé en défaut}, que ce changement soit
\textbf{tracé} (qui, quand, avec quel retard), que le service de recouvrement ait été \textbf{alerté},
et que la provision de fin de mois en tienne compte. Et pendant ce temps, l'application d'encaissement
doit continuer à enregistrer les paiements des autres clients, sans jamais pouvoir effacer un paiement
ni modifier un contrat.
\end{situation}

% =============================================================================
\section{Cours et notions importantes}

\subsection{Où placer la logique métier ?}

Une règle métier peut vivre dans l'application (Python, Java\dots) ou dans la base. Les deux choix se
défendent ; en entretien, on attend que tu saches les comparer.

\begin{center}\small
\begin{tabularx}{\linewidth}{@{}B{3.4cm} L L@{}}
\toprule
& Dans la base (fonctions, triggers) & Dans l'application \\
\midrule
Contournement & impossible : tous les accès passent par la base & une requête directe ou un script
contourne la règle \\
Cohérence & une seule implémentation & une par application qui écrit dans la base \\
Performance & au plus près des données, ensembliste & allers-retours réseau \\
Tests, versions & plus difficile (outils dédiés) & outils habituels du développeur \\
Lisibilité & logique « cachée » dans la base & visible dans le code \\
\bottomrule
\end{tabularx}
\end{center}

\begin{regle}[La règle pratique]
On met dans la base ce qui \textbf{ne doit jamais être contourné} : l'intégrité des données, la
traçabilité, les règles réglementaires simples (le passage en défaut). On laisse dans l'application
l'interface, l'orchestration et les calculs complexes (l'entraînement d'un modèle de ML, par exemple).
\end{regle}

\subsection{Les fonctions : SQL ou PL/pgSQL}

Une fonction enregistre un calcul sous un nom, avec des paramètres. PostgreSQL en propose deux
langages principaux :
\begin{itemize}
\item \textbf{SQL} : le corps est une simple requête. Idéal pour une formule.
\item \textbf{PL/pgSQL} : un vrai langage procédural, avec variables, conditions, boucles et gestion des
erreurs.
\end{itemize}

\begin{lstlisting}[style=sql]
CREATE OR REPLACE FUNCTION fn_date_reference() RETURNS DATE
LANGUAGE sql STABLE AS $$
    SELECT valeur_date FROM parametres_risque WHERE cle = 'date_reference'
$$;
\end{lstlisting}

Le corps est encadré par \code{$$ ... $$} (\emph{dollar quoting}) : tout ce qui est entre les deux est
du texte, sans avoir à doubler les apostrophes. Un bloc PL/pgSQL a toujours la même structure :

\begin{lstlisting}[style=sql]
CREATE OR REPLACE FUNCTION fn_parametre(p_cle TEXT) RETURNS NUMERIC
LANGUAGE plpgsql STABLE AS $$
DECLARE                                   -- 1. les variables
    v_valeur NUMERIC;
BEGIN                                     -- 2. les instructions
    SELECT valeur INTO STRICT v_valeur FROM parametres_risque WHERE cle = p_cle;
    RETURN v_valeur;
EXCEPTION WHEN NO_DATA_FOUND THEN         -- 3. la gestion des erreurs (facultative)
    RAISE EXCEPTION 'Paramètre inconnu : %', p_cle;
END $$;
\end{lstlisting}

\code{INTO STRICT} exige exactement une ligne : zéro ligne lève \code{NO_DATA_FOUND}, plusieurs lignes
\code{TOO_MANY_ROWS}. Sans \code{STRICT}, une clé mal orthographiée donnerait silencieusement NULL, et
tout calcul qui en dépend deviendrait NULL sans que personne ne s'en aperçoive. \code{RAISE EXCEPTION}
annule l'instruction en cours avec un message ; les \code{\%} sont remplacés par les valeurs qui
suivent.

\subsubsection{La volatilité : promettre quelque chose à l'optimiseur}

\begin{center}\small
\begin{tabularx}{\linewidth}{@{}P{2.4cm} L P{4.8cm}@{}}
\toprule
\textbf{Catégorie} & \textbf{Promesse} & \textbf{Dans le projet} \\
\midrule
\code{IMMUTABLE} & mêmes arguments $\Rightarrow$ même résultat, pour toujours ; ne lit pas la base &
\code{fn_pd_ifrs9} (une formule) \\
\code{STABLE} & même résultat pendant toute une instruction ; peut lire la base & \code{fn_parametre},
\code{fn_jours_retard} \\
\code{VOLATILE} & aucune promesse (par défaut) ; peut modifier la base & \code{random()}, les fonctions
de trigger \\
\bottomrule
\end{tabularx}
\end{center}

Ces promesses permettent à PostgreSQL de calculer une fois au lieu de mille. Mentir (déclarer
\code{IMMUTABLE} une fonction qui lit une table) peut donner des résultats faux, par exemple dans un
index.

\subsubsection{Ensembliste ou ligne à ligne}

Le retard d'\textbf{un} prêt s'écrit naturellement comme une fonction, \code{fn_jours_retard(pret, date)}.
Mais l'appeler pour 32\,410 prêts, c'est exécuter 32\,410 requêtes. Grâce à l'index partiel, chacune
prend quelques microsecondes : 0,25 seconde au total. Sans l'index, chacune relirait 1,1 million de
lignes : plus de 20 minutes. La version \textbf{ensembliste}, \code{fn_situation_prets(date)}, calcule
tout le portefeuille en une seule requête et renvoie une table (\code{RETURNS TABLE}).

\begin{piege}{Piège --- La fonction SQL « dépliée »}
PostgreSQL recopie le corps d'une fonction SQL simple dans la requête qui l'appelle (\emph{inlining}).
C'est souvent un gain, mais si l'argument est lui-même un appel de fonction, comme
\code{fn_situation_prets(fn_date_reference())}, cet argument peut se retrouver recalculé pour chacune
des 1,1 million de lignes parcourues : 3,5 secondes au lieu de 0,2. Écrite en PL/pgSQL, la fonction
reçoit son argument déjà calculé, une seule fois.
\end{piege}

\subsection{Les triggers}

\begin{analogie}{le détecteur de fumée}
Un détecteur de fumée ne fait rien tant qu'il n'y a pas de fumée. Il n'a pas besoin qu'on pense à lui :
dès que l'événement se produit, il réagit, toujours de la même façon. Un trigger est un détecteur
installé sur une table.
\end{analogie}

\begin{definition}{Trigger (déclencheur)}
Une fonction que PostgreSQL exécute \textbf{automatiquement} quand un événement (\code{INSERT},
\code{UPDATE}, \code{DELETE}, \code{TRUNCATE}) touche une table. La fonction est déclarée
\code{RETURNS trigger} ; le trigger l'attache à la table et précise quand l'appeler.
\end{definition}

\begin{lstlisting}[style=sql]
CREATE TRIGGER trg_journal_statut
    AFTER UPDATE OF statut_pret ON demandes_pret        -- quand : après la modification de cette colonne
    FOR EACH ROW                                        -- pour chaque ligne modifiée
    WHEN (OLD.statut_pret IS DISTINCT FROM NEW.statut_pret)   -- seulement si le statut change vraiment
    EXECUTE FUNCTION fn_trg_journal_statut();           -- quoi faire
\end{lstlisting}

\begin{center}\small
\begin{tabularx}{\linewidth}{@{}P{2.9cm} P{4.2cm} L@{}}
\toprule
\textbf{Choix} & \textbf{Options} & \textbf{Ce que ça change} \\
\midrule
Moment & \code{BEFORE}, \code{AFTER} & \code{BEFORE} peut refuser ou modifier la ligne avant écriture ;
\code{AFTER} voit le résultat final et réagit \\
Événement & \code{INSERT}, \code{UPDATE [OF col]}, \code{DELETE} & \code{UPDATE OF} limite le trigger à
certaines colonnes \\
Niveau & \code{FOR EACH ROW}, \code{FOR EACH STATEMENT} & une fois par ligne touchée, ou une fois par
instruction \\
Condition & \code{WHEN (...)} & la fonction n'est même pas appelée si la condition est fausse \\
\bottomrule
\end{tabularx}
\end{center}

Dans une fonction de trigger par ligne, deux variables décrivent la ligne : \code{OLD} (avant) et
\code{NEW} (après). La valeur renvoyée a un sens précis :
\begin{itemize}
\item dans un trigger \code{BEFORE} par ligne, \code{RETURN NEW} fait écrire la ligne \code{NEW}
(éventuellement modifiée par la fonction) ; \code{RETURN NULL} annule silencieusement l'opération sur
cette ligne ; \code{RAISE EXCEPTION} annule toute l'instruction avec une erreur ;
\item dans un trigger \code{AFTER}, la valeur renvoyée est ignorée : on écrit \code{RETURN NULL}.
\end{itemize}

Un trigger \code{AFTER ... FOR EACH STATEMENT} peut recevoir \textbf{toutes} les lignes modifiées par
l'instruction, sous forme de tables : les \textbf{tables de transition}.
\begin{lstlisting}[style=sql]
AFTER UPDATE ON echeancier_remboursements
REFERENCING OLD TABLE AS anciennes NEW TABLE AS nouvelles   -- toutes les lignes avant / après
FOR EACH STATEMENT EXECUTE FUNCTION ...
\end{lstlisting}
Une instruction qui enregistre 1\,000 paiements déclenche alors \textbf{un} recalcul ensembliste au
lieu de 1\,000.

\begin{figure}[ht]
\centering
\begin{tikzpicture}[font=\footnotesize, >=Stealth,
  b/.style={draw, rounded corners=2pt, minimum height=0.75cm, align=center, inner sep=3pt},
  av/.style={b, draw=orangeEx, fill=orangeFond}, ap/.style={b, draw=bleuML, fill=bleuFond},
  ct/.style={b, draw=grisTexte, fill=grisFond}, st/.style={b, draw=vertCadre, fill=vertFond}]
  \node[font=\small\bfseries, anchor=west] at (-0.4,0.75) {UPDATE de 2 lignes :};
  \node[av] (a1) at (1,0) {BEFORE\\ligne 1};
  \node[ct] (c1) at (2.75,0) {contraintes\\ligne 1};
  \node[av] (a2) at (4.5,0) {BEFORE\\ligne 2};
  \node[ct] (c2) at (6.25,0) {contraintes\\ligne 2};
  \node[ap] (p1) at (8.15,0) {AFTER\\ligne 1};
  \node[ap] (p2) at (9.9,0) {AFTER\\ligne 2};
  \node[st] (s) at (12.1,0) {AFTER\\STATEMENT};
  \foreach \x/\y in {a1/c1,c1/a2,a2/c2,c2/p1,p1/p2,p2/s} \draw[->, thick, grisTexte] (\x) -- (\y);
  \draw[decorate, decoration={brace, amplitude=4pt, mirror}] (0.3,-0.55) -- (7.0,-0.55)
     node[midway, below=5pt] {pendant l'écriture, ligne par ligne};
  \draw[decorate, decoration={brace, amplitude=4pt, mirror}] (7.45,-0.55) -- (13.0,-0.55)
     node[midway, below=5pt] {à la fin de l'instruction};
\end{tikzpicture}
\caption{L'ordre d'exécution des triggers pour une instruction qui modifie deux lignes. Si plusieurs
triggers ont le même moment, ils s'exécutent dans l'ordre alphabétique de leur nom.}
\label{fig:triggers}
\end{figure}

\begin{piege}{Piège --- Les triggers en cascade}
Un trigger qui modifie une autre table peut en déclencher d'autres : dans ton projet, un paiement
déclenche le recalcul du statut, qui déclenche le journal. C'est voulu. Mais un trigger qui modifie
\textbf{sa propre table} se redéclenche lui-même, à l'infini. Et une erreur dans n'importe quel trigger
annule toute l'instruction d'origine. Un trigger est aussi \textbf{invisible} pour qui lit une requête :
documente-les.
\end{piege}

\subsection{Les procédures stockées}

\begin{center}\small
\begin{tabularx}{\linewidth}{@{}B{3.6cm} L L@{}}
\toprule
& \code{FUNCTION} & \code{PROCEDURE} \\
\midrule
Appel & dans une requête : \code{SELECT f(x)} & instruction seule : \code{CALL p(x)} \\
Résultat & renvoie une valeur ou une table & ne renvoie rien (ou des paramètres \code{INOUT}) \\
Transactions & ne peut pas valider & peut faire \code{COMMIT} entre deux étapes (si elle n'est pas
appelée dans un bloc \code{BEGIN}) \\
Usage typique & un calcul & une opération métier : encaisser, clôturer \\
\bottomrule
\end{tabularx}
\end{center}

Une procédure PL/pgSQL utilise les mêmes briques qu'une fonction, plus quelques outils fréquents :
\begin{lstlisting}[style=sql]
FOR ech IN SELECT ... ORDER BY num_echeance LOOP   -- parcourir le résultat d'une requête
    EXIT WHEN v_restant = 0;                       -- sortir de la boucle
    ...
END LOOP;
GET DIAGNOSTICS v_nb = ROW_COUNT;                  -- nombre de lignes touchées par la dernière instruction
RAISE NOTICE '% statut(s) mis à jour', v_nb;       -- un message d'information (pas une erreur)
\end{lstlisting}

\subsubsection{Deux paiements en même temps : le verrou}

\begin{imagine}
Le prêt 61 a une échéance impayée de 704,41~€. Au même instant, deux virements arrivent : 704,41~€ par
l'application mobile, et 704,41~€ par le guichet. Les deux procédures lisent « échéance 18 impayée »,
et les deux l'affectent à l'échéance 18. Résultat : un paiement est perdu, ou la contrainte
\code{paiement_borne} fait échouer l'un des deux au hasard.
\end{imagine}

La solution est de \textbf{verrouiller le prêt} au début de la procédure :
\code{SELECT ... FROM demandes_pret WHERE pret_id = p FOR UPDATE}. Le second paiement attend que le
premier soit validé, puis voit l'échéance 18 déjà payée et passe à la 19. Les paiements du même prêt
sont traités l'un après l'autre ; ceux des autres prêts continuent en parallèle.

\subsection{Le temps qui passe : le traitement de nuit}

Comme on l'a vu dans la partie 2, \textbf{le passage du temps n'est pas un événement}. Le retard d'Awa
franchit 90 jours un dimanche sans qu'aucune ligne ne change : aucun trigger ne peut le détecter. Il faut
une procédure, \code{actualiser_statuts()}, qui recalcule les statuts de tout le portefeuille, et un
\textbf{ordonnanceur} qui la lance chaque nuit. Sur un serveur, c'est une ligne de \code{cron} :
\begin{lstlisting}[style=shell]
0 2 * * *  psql -d risk_analytics_db -c "CALL actualiser_statuts();"
\end{lstlisting}
ou l'extension \code{pg_cron}, qui planifie directement dans PostgreSQL. Les deux mécanismes partagent
la même règle de statut, écrite une seule fois dans une fonction.

\subsection{Rejouer sans dupliquer : l'idempotence}

Une clôture mensuelle peut échouer à mi-parcours, ou devoir être relancée après une correction. Elle
doit donner le même résultat qu'on la lance une ou trois fois : elle doit être \textbf{idempotente}.
L'instruction \code{INSERT ... ON CONFLICT} (l'« upsert ») le permet :
\begin{lstlisting}[style=sql]
INSERT INTO provisions_mensuelles (date_arrete, pret_id, ...) SELECT ...
ON CONFLICT (date_arrete, pret_id) DO UPDATE             -- la ligne existe déjà : on la remplace
    SET ecl = EXCLUDED.ecl, ...;                          -- EXCLUDED = la ligne qu'on voulait insérer
\end{lstlisting}

\subsection{Alerter : LISTEN et NOTIFY}

PostgreSQL peut envoyer des messages aux programmes connectés. Un programme s'abonne à un canal avec
\code{LISTEN alertes_defaut} ; une fonction envoie un message avec \code{pg_notify('alertes_defaut',
'...')}. Le message n'est délivré qu'au \textbf{moment du COMMIT} : si la transaction est annulée,
aucune fausse alerte n'est envoyée. Voici ce que psql affiche quand le prêt d'Awa passe en défaut :
\begin{lstlisting}[style=sortie]
Asynchronous notification "alertes_defaut" with payload "{"pret_id" : 32582, "jours_retard" : 98}"
received from server process with PID 87868.
\end{lstlisting}
Dans une application Python, la boucle \code{for n in conn.notifies(): ...} de psycopg reçoit ces
messages et peut, par exemple, créer une tâche pour le service de recouvrement.

\subsection{La sécurité : le moindre privilège}

Chaque utilisateur ne doit avoir que les droits dont il a besoin. PostgreSQL gère des \textbf{rôles} et
des droits (\code{GRANT}, \code{REVOKE}) :
\begin{itemize}
\item l'analyste risque \textbf{lit} tout, ne modifie rien : \code{GRANT SELECT} ;
\item l'application d'encaissement n'écrit \textbf{jamais} directement dans les tables : elle n'a le droit
que d'appeler \code{enregistrer_paiement}. La procédure est déclarée \code{SECURITY DEFINER} : elle
s'exécute avec les droits de son propriétaire, et fait donc ce que l'application ne peut pas faire
elle-même, mais seulement de la façon prévue.
\end{itemize}

\begin{piege}{Piège --- SECURITY DEFINER sans search\_path}
Une fonction \code{SECURITY DEFINER} cherche les tables et fonctions qu'elle utilise dans le
\code{search_path} de l'utilisateur qui l'appelle. Un utilisateur malveillant pourrait créer une fausse
table homonyme dans son propre schéma et détourner la procédure, qui s'exécute avec des droits élevés.
Il faut toujours figer le chemin : \code{SET search_path = public, pg_temp}.
\end{piege}

\subsection{La règle du statut}

Toute l'automatisation repose sur une règle, évaluée \textbf{dans cet ordre} (la première condition
vraie l'emporte) :

\begin{center}\small
\begin{tabular}{@{}clll@{}}
\toprule
Ordre & Condition & Nouveau statut & Justification \\
\midrule
1 & reste dû $= 0$ & \code{solde} & tout est remboursé \\
2 & statut actuel $=$ en défaut & \code{en_defaut} & pas de sortie automatique (période probatoire) \\
3 & DPD $> 90$ & \code{en_defaut} & article 178 du CRR \\
4 & DPD $> 0$ & \code{en_retard} & au moins une échéance impayée \\
5 & sinon & \code{sain} & \\
\bottomrule
\end{tabular}
\end{center}

% =============================================================================
\section{Exercices guidés}

\begin{exercice}{3.1}{Des fonctions pour ne jamais réécrire une règle métier}
Écris les fonctions suivantes, puis teste-les :
\begin{enumerate}[label=\alph*)]
\item \code{fn_parametre(cle)} (erreur explicite si la clé n'existe pas) et \code{fn_date_reference()} ;
\item \code{fn_jours_retard(pret, date)} : le retard d'un prêt à une date quelconque, passée ou
actuelle ;
\item \code{fn_situation_prets(date)} : pour tout le portefeuille à cette date, le retard, l'encours, le
reste dû et les mois restants ;
\item \code{fn_stage_ifrs9(jours_retard)} et \code{fn_pd_ifrs9(stage, pd, mois_restants, duree)}.
\end{enumerate}
Tests : le retard du prêt 3 ; la concordance entre \code{fn_jours_retard} et \code{fn_situation_prets} ;
la concordance entre \code{fn_situation_prets} et la vue du chapitre 2.
\end{exercice}

\begin{correction}{3.1}
\textbf{a) et b)} Le retard « à une date quelconque » demande une définition plus générale que celle de
la vue : une échéance payée \emph{après} la date $D$ était impayée \emph{à} la date $D$. D'où la
condition \code{date_paiement IS NULL OR date_paiement > p_date}. C'est elle qui permettra, au chapitre
suivant, de reconstituer le passé (provisions des mois précédents, matrice de transition).

\begin{lstlisting}[style=long]
-- a) Lire un paramètre : une erreur explicite vaut mieux qu'un NULL silencieux
CREATE OR REPLACE FUNCTION fn_parametre(p_cle TEXT) RETURNS NUMERIC
LANGUAGE plpgsql STABLE AS $$
DECLARE
    v_valeur NUMERIC;
BEGIN
    SELECT valeur INTO STRICT v_valeur FROM parametres_risque WHERE cle = p_cle;
    RETURN v_valeur;
EXCEPTION WHEN NO_DATA_FOUND THEN
    RAISE EXCEPTION 'Paramètre inconnu : %', p_cle;
END $$;

CREATE OR REPLACE FUNCTION fn_date_reference() RETURNS DATE
LANGUAGE sql STABLE AS $$
    SELECT valeur_date FROM parametres_risque WHERE cle = 'date_reference'
$$;

-- b) Le retard d'UN prêt à une date quelconque (passée ou actuelle)
CREATE OR REPLACE FUNCTION fn_jours_retard(p_pret_id BIGINT, p_date DATE) RETURNS INTEGER
LANGUAGE sql STABLE AS $$
    SELECT coalesce(p_date - min(e.date_echeance), 0)
    FROM echeancier_remboursements e
    WHERE e.pret_id = p_pret_id
      AND e.date_echeance <= p_date
      AND (e.date_paiement IS NULL OR e.date_paiement > p_date)    -- pas encore soldée à cette date
$$;
\end{lstlisting}

\textbf{c)} La version ensembliste, en PL/pgSQL pour éviter le piège du dépliage. \code{RETURNS TABLE}
déclare les colonnes de sortie ; attention, leurs noms deviennent des variables dans la fonction : on
préfixe donc toutes les colonnes des tables (\code{e.pret_id}) pour éviter toute ambiguïté.

\begin{lstlisting}[style=long]
CREATE OR REPLACE FUNCTION fn_situation_prets(p_date DATE)
RETURNS TABLE (pret_id BIGINT, jours_retard INTEGER, encours NUMERIC, reste_du NUMERIC, mois_restants INTEGER)
LANGUAGE plpgsql STABLE AS $$
BEGIN
    RETURN QUERY
    SELECT e.pret_id,
           coalesce(p_date - min(e.date_echeance)
                      FILTER (WHERE e.date_echeance <= p_date
                                AND (e.date_paiement IS NULL OR e.date_paiement > p_date)), 0),
           -- capital non remboursé à la date : impayé partiel au prorata, soldé après la date = encore dû
           round(sum(e.part_capital * CASE WHEN e.date_paiement IS NULL THEN 1 - e.montant_paye / e.montant_du
                                           WHEN e.date_paiement > p_date THEN 1
                                           ELSE 0 END), 2),
           sum(e.montant_du - e.montant_paye),
           (count(*) FILTER (WHERE e.date_echeance > p_date))::int
    FROM echeancier_remboursements e
    JOIN demandes_pret d ON d.pret_id = e.pret_id
    WHERE d.statut_dossier = 'accorde' AND d.date_decision <= p_date
    GROUP BY e.pret_id;                    -- colonnes préfixées (e., d.) : pas de conflit avec les noms en sortie
END $$;
\end{lstlisting}

\textbf{d)} Les règles IFRS~9 du chapitre 2, enfin nommées et réutilisables :
\begin{lstlisting}[style=long]
CREATE OR REPLACE FUNCTION fn_stage_ifrs9(p_jours_retard INTEGER) RETURNS SMALLINT
LANGUAGE sql STABLE AS $$
    SELECT (CASE WHEN p_jours_retard > fn_parametre('seuil_defaut_jours') THEN 3
                 WHEN p_jours_retard > fn_parametre('seuil_stage2_jours') THEN 2
                 ELSE 1 END)::smallint
$$;

CREATE OR REPLACE FUNCTION fn_pd_ifrs9(p_stage SMALLINT, p_pd_vie NUMERIC, p_mois_restants INTEGER, p_duree INTEGER)
RETURNS NUMERIC LANGUAGE sql IMMUTABLE AS $$
    SELECT CASE p_stage
               WHEN 3 THEN 1
               WHEN 2 THEN 1 - (1 - p_pd_vie) ^ (greatest(p_mois_restants, 1)::numeric / p_duree)
               ELSE        1 - (1 - p_pd_vie) ^ (least(greatest(p_mois_restants, 1), 12)::numeric / p_duree)
           END
$$;
\end{lstlisting}
\code{fn_stage_ifrs9} est \code{STABLE} (elle lit les paramètres), \code{fn_pd_ifrs9} est
\code{IMMUTABLE} (une pure formule).

\textbf{Les tests.} Une fonction non testée n'existe pas. On vérifie que les deux versions du calcul
donnent la même chose, et qu'elles concordent avec la vue du chapitre 2 :
\begin{lstlisting}[style=long]
SELECT fn_jours_retard(3, fn_date_reference()) AS retard_pret_3;

SELECT count(*) AS desaccords_fonction_vs_ensembliste          -- attendu : 0
FROM fn_situation_prets(fn_date_reference()) s
WHERE s.jours_retard <> fn_jours_retard(s.pret_id, fn_date_reference());

SELECT count(*) AS desaccords_vue_vs_fonction                  -- attendu : 0
FROM v_situation_prets v
JOIN fn_situation_prets(fn_date_reference()) s USING (pret_id)
WHERE (v.jours_retard, v.encours) IS DISTINCT FROM (s.jours_retard, s.encours);
\end{lstlisting}

\begin{resultat}
Retard du prêt 3 : \textbf{1\,170 jours}. Désaccords : \textbf{0} et \textbf{0}.
\end{resultat}
La comparaison \code{(a, b) IS DISTINCT FROM (c, d)} compare deux couples d'un coup, en traitant
correctement les NULL.
\end{correction}

\begin{exercice}{3.2}{Trois triggers : protéger, recalculer, tracer}
\begin{enumerate}[label=\alph*)]
\item \textbf{Protéger} l'échéancier : un \code{UPDATE} ne peut pas modifier les termes du contrat
(numéro, date, montants, CRD) ni diminuer un montant payé. Une échéance qui devient entièrement payée
sans date de paiement reçoit la date de référence.
\item Écrire la \textbf{règle du statut} dans une fonction \code{fn_regle_statut(statut_actuel,
jours_retard, reste_du)}.
\item \textbf{Recalculer} : après chaque instruction qui modifie des montants payés, recalculer le
statut des prêts concernés, en un seul passage.
\item \textbf{Tracer} : chaque changement de \code{statut_pret} est inscrit dans \code{journal_statuts},
et un passage en défaut envoie une notification sur le canal \code{alertes_defaut}.
\end{enumerate}
\end{exercice}

\begin{correction}{3.2}
\textbf{a) Le gardien : BEFORE UPDATE, par ligne.} Seul un trigger \code{BEFORE} peut refuser une
modification ou compléter la ligne avant son écriture. La comparaison de deux couples par
\code{IS DISTINCT FROM} détecte n'importe quel changement des termes du contrat, y compris vers NULL.

\begin{lstlisting}[style=long]
CREATE OR REPLACE FUNCTION fn_trg_proteger_echeance() RETURNS trigger
LANGUAGE plpgsql AS $$
BEGIN
    IF (NEW.pret_id, NEW.num_echeance, NEW.date_echeance, NEW.montant_du,
        NEW.part_capital, NEW.part_interets, NEW.capital_restant_du)
       IS DISTINCT FROM
       (OLD.pret_id, OLD.num_echeance, OLD.date_echeance, OLD.montant_du,
        OLD.part_capital, OLD.part_interets, OLD.capital_restant_du) THEN
        RAISE EXCEPTION 'Échéance % : les termes du contrat ne se modifient pas', OLD.echeance_id
            USING HINT = 'Un changement de contrat passe par un avenant et un nouvel échéancier.';
    END IF;
    IF NEW.montant_paye < OLD.montant_paye THEN
        RAISE EXCEPTION 'Échéance % : un paiement ne s''efface pas (% → %)',
                        OLD.echeance_id, OLD.montant_paye, NEW.montant_paye
            USING HINT = 'Enregistrer une écriture de régularisation.';
    END IF;
    IF NEW.montant_paye = NEW.montant_du AND NEW.date_paiement IS NULL THEN
        NEW.date_paiement := fn_date_reference();      -- échéance soldée : la date est complétée
    END IF;
    RETURN NEW;                                        -- BEFORE : la ligne renvoyée est celle qui sera écrite
END $$;

DROP TRIGGER IF EXISTS trg_proteger_echeance ON echeancier_remboursements;
CREATE TRIGGER trg_proteger_echeance
    BEFORE UPDATE ON echeancier_remboursements
    FOR EACH ROW EXECUTE FUNCTION fn_trg_proteger_echeance();
\end{lstlisting}

Test, avec deux instructions qui doivent échouer :
\begin{lstlisting}[style=sql]
UPDATE echeancier_remboursements SET montant_du = 1   WHERE echeance_id = 1;
UPDATE echeancier_remboursements SET montant_paye = 0 WHERE echeance_id = 1;
\end{lstlisting}
\begin{lstlisting}[style=sortie]
ERROR:  Échéance 1 : les termes du contrat ne se modifient pas
HINT:  Un changement de contrat passe par un avenant et un nouvel échéancier.
ERROR:  Échéance 1 : un paiement ne s'efface pas (88.45 → 0.00)
HINT:  Enregistrer une écriture de régularisation.
\end{lstlisting}

\textbf{b) La règle, écrite une fois.} Elle sera utilisée par le trigger \emph{et} par le traitement de
nuit : impossible qu'ils divergent.
\begin{lstlisting}[style=long]
CREATE OR REPLACE FUNCTION fn_regle_statut(p_statut_actuel VARCHAR, p_jours_retard INTEGER, p_reste_du NUMERIC)
RETURNS VARCHAR LANGUAGE sql STABLE AS $$
    SELECT CASE
               WHEN p_reste_du = 0                                     THEN 'solde'
               -- pas de sortie automatique du défaut : période probatoire et décision humaine
               WHEN p_statut_actuel = 'en_defaut'                      THEN 'en_defaut'
               WHEN p_jours_retard > fn_parametre('seuil_defaut_jours') THEN 'en_defaut'
               WHEN p_jours_retard > 0                                 THEN 'en_retard'
               ELSE                                                         'sain'
           END
$$;
\end{lstlisting}

\textbf{c) Le recalcul : AFTER UPDATE, par instruction, avec tables de transition.} La fonction joint
\code{anciennes} et \code{nouvelles} pour ne garder que les prêts dont un montant payé a changé, puis
met à jour leur statut. La condition \code{IS DISTINCT FROM} évite d'écrire (et donc de journaliser)
un statut inchangé.
\begin{lstlisting}[style=long]
CREATE OR REPLACE FUNCTION fn_trg_statut_apres_paiement() RETURNS trigger
LANGUAGE plpgsql AS $$
BEGIN
    WITH prets_payes AS (
        SELECT DISTINCT n.pret_id
        FROM nouvelles n
        JOIN anciennes a USING (echeance_id)
        WHERE n.montant_paye <> a.montant_paye
    ),
    recalcul AS (
        SELECT d.pret_id,
               fn_regle_statut(d.statut_pret,
                               fn_jours_retard(d.pret_id, fn_date_reference()),
                               (SELECT sum(e.montant_du - e.montant_paye)
                                FROM echeancier_remboursements e WHERE e.pret_id = d.pret_id)) AS statut
        FROM demandes_pret d
        JOIN prets_payes USING (pret_id)
    )
    UPDATE demandes_pret d
    SET statut_pret = r.statut
    FROM recalcul r
    WHERE d.pret_id = r.pret_id
      AND d.statut_pret IS DISTINCT FROM r.statut;      -- n'écrit (et ne journalise) que les vrais changements
    RETURN NULL;                                        -- AFTER : la valeur renvoyée est ignorée
END $$;

DROP TRIGGER IF EXISTS trg_statut_apres_paiement ON echeancier_remboursements;
CREATE TRIGGER trg_statut_apres_paiement
    AFTER UPDATE ON echeancier_remboursements
    REFERENCING OLD TABLE AS anciennes NEW TABLE AS nouvelles
    FOR EACH STATEMENT EXECUTE FUNCTION fn_trg_statut_apres_paiement();
\end{lstlisting}

\textbf{d) Le journal et l'alerte : AFTER UPDATE OF, par ligne, avec WHEN.} \code{session_user} et
\code{now()} sont remplis par les valeurs par défaut de la table.
\begin{lstlisting}[style=long]
CREATE OR REPLACE FUNCTION fn_trg_journal_statut() RETURNS trigger
LANGUAGE plpgsql AS $$
DECLARE
    v_retard INTEGER := fn_jours_retard(NEW.pret_id, fn_date_reference());
BEGIN
    INSERT INTO journal_statuts (pret_id, ancien_statut, nouveau_statut, jours_retard, date_effet)
    VALUES (NEW.pret_id, OLD.statut_pret, NEW.statut_pret, v_retard, fn_date_reference());
    IF NEW.statut_pret = 'en_defaut' THEN
        -- alerte : tout programme abonné (LISTEN alertes_defaut) la reçoit quand la transaction est validée
        PERFORM pg_notify('alertes_defaut',
                          json_build_object('pret_id', NEW.pret_id, 'jours_retard', v_retard)::text);
    END IF;
    RETURN NULL;
END $$;

DROP TRIGGER IF EXISTS trg_journal_statut ON demandes_pret;
CREATE TRIGGER trg_journal_statut
    AFTER UPDATE OF statut_pret ON demandes_pret
    FOR EACH ROW
    WHEN (OLD.statut_pret IS DISTINCT FROM NEW.statut_pret)    -- la fonction n'est même pas appelée sinon
    EXECUTE FUNCTION fn_trg_journal_statut();
\end{lstlisting}

\textbf{La chaîne complète.} Un paiement modifie \code{echeancier_remboursements} $\rightarrow$ le gardien
le valide $\rightarrow$ le recalcul met à jour \code{demandes_pret} $\rightarrow$ le journal enregistre le
changement et, en cas de défaut, alerte. Trois triggers, trois responsabilités.
\end{correction}

\begin{exercice}{3.3}{Deux procédures : l'encaissement et le traitement de nuit}
\begin{enumerate}[label=\alph*)]
\item \code{enregistrer_paiement(pret, montant, date)} : vérifie le montant et le prêt, verrouille le
prêt, refuse un paiement supérieur au reste dû, puis affecte le paiement aux échéances impayées
\textbf{de la plus ancienne à la plus récente}.
\item \code{actualiser_statuts(pret)} : recalcule le statut d'un prêt, ou de tout le portefeuille si
aucun prêt n'est donné. Lance-la une première fois : tous les prêts chargés sont « sain ».
\item Dans une transaction annulée à la fin, joue trois scénarios : le prêt 67 règle tout son impayé
(916,55~€) ; le prêt 61 paie 300~€ ; le prêt 9 règle tous ses impayés (9\,422,92~€). Observe les
statuts et le journal.
\end{enumerate}
\end{exercice}

\begin{correction}{3.3}
\textbf{a) L'encaissement.} Les contrôles d'abord, le verrou ensuite, puis une boucle sur les échéances
impayées dans l'ordre. Chaque \code{UPDATE} passe par le gardien ; le recalcul du statut se fait
automatiquement, après chaque instruction.
\begin{lstlisting}[style=long]
CREATE OR REPLACE PROCEDURE enregistrer_paiement(p_pret_id BIGINT, p_montant NUMERIC, p_date DATE DEFAULT NULL)
LANGUAGE plpgsql AS $$
DECLARE
    v_date     DATE    := coalesce(p_date, fn_date_reference());
    v_restant  NUMERIC := p_montant;
    v_dossier  VARCHAR;
    v_du_total NUMERIC;
    v_part     NUMERIC;
    ech        RECORD;
BEGIN
    IF p_montant IS NULL OR p_montant <= 0 THEN
        RAISE EXCEPTION 'Montant invalide : %', p_montant;
    END IF;
    -- verrou sur le prêt : deux paiements simultanés du même prêt passent l'un après l'autre
    SELECT statut_dossier INTO v_dossier FROM demandes_pret WHERE pret_id = p_pret_id FOR UPDATE;
    IF NOT FOUND THEN
        RAISE EXCEPTION 'Prêt % inconnu', p_pret_id;
    ELSIF v_dossier <> 'accorde' THEN
        RAISE EXCEPTION 'Le prêt % n''est pas accordé', p_pret_id;
    END IF;
    SELECT sum(montant_du - montant_paye) INTO v_du_total
    FROM echeancier_remboursements WHERE pret_id = p_pret_id;
    IF p_montant > v_du_total THEN
        RAISE EXCEPTION 'Paiement de % € supérieur au reste dû du prêt % (% €)', p_montant, p_pret_id, v_du_total;
    END IF;

    FOR ech IN
        SELECT echeance_id, montant_du - montant_paye AS manquant
        FROM echeancier_remboursements
        WHERE pret_id = p_pret_id AND montant_paye < montant_du
        ORDER BY num_echeance                          -- la plus ancienne d'abord
    LOOP
        EXIT WHEN v_restant = 0;
        v_part := least(v_restant, ech.manquant);
        UPDATE echeancier_remboursements
        SET montant_paye  = montant_paye + v_part,
            date_paiement = CASE WHEN montant_paye + v_part = montant_du THEN v_date END
        WHERE echeance_id = ech.echeance_id;
        v_restant := v_restant - v_part;
    END LOOP;
    RAISE NOTICE 'Prêt % : % € encaissés le %', p_pret_id, p_montant, v_date;
END $$;
\end{lstlisting}

\textbf{b) Le traitement de nuit.} Une seule instruction ensembliste, qui réutilise les deux
fonctions de l'exercice précédent : \code{fn_situation_prets} et \code{fn_regle_statut}. Grâce au
paramètre \code{DEFAULT NULL}, la même procédure traite un prêt ou tout le portefeuille.
\begin{lstlisting}[style=long]
CREATE OR REPLACE PROCEDURE actualiser_statuts(p_pret_id BIGINT DEFAULT NULL)   -- NULL = tout le portefeuille
LANGUAGE plpgsql AS $$
DECLARE
    v_nb INTEGER;
BEGIN
    WITH recalcul AS (
        SELECT d.pret_id, fn_regle_statut(d.statut_pret, s.jours_retard, s.reste_du) AS statut
        FROM demandes_pret d
        JOIN fn_situation_prets(fn_date_reference()) s USING (pret_id)
        WHERE p_pret_id IS NULL OR d.pret_id = p_pret_id
    )
    UPDATE demandes_pret d
    SET statut_pret = r.statut
    FROM recalcul r
    WHERE d.pret_id = r.pret_id AND d.statut_pret IS DISTINCT FROM r.statut;
    GET DIAGNOSTICS v_nb = ROW_COUNT;
    RAISE NOTICE '% statut(s) mis à jour au %', v_nb, fn_date_reference();
END $$;

-- Premier passage : tous les prêts chargés sont « sain », le traitement leur donne leur vrai statut
CALL actualiser_statuts();
SELECT statut_pret, count(*) AS nb_prets FROM demandes_pret GROUP BY statut_pret ORDER BY statut_pret;
\end{lstlisting}

\begin{resultat}
\code{NOTICE: 18930 statut(s) mis à jour au 2026-09-30} : 7\,088 \code{en_defaut}, 933 \code{en_retard},
13\,480 \code{sain}, 10\,909 \code{solde}. Le journal contient désormais 18\,930 lignes.
\end{resultat}

\textbf{c) Les scénarios.} \code{BEGIN ... ROLLBACK} : on observe tout, puis la base revient à son état
initial.
\begin{lstlisting}[style=long]
BEGIN;
CALL enregistrer_paiement(67, 916.55);     -- 1. prêt en retard de 3 jours : il règle tout son impayé
CALL enregistrer_paiement(61, 300);        -- 2. prêt en retard de 46 jours : paiement partiel
CALL enregistrer_paiement(9, 9422.92);     -- 3. prêt en défaut (21 mois restants) : il règle tous ses impayés
SELECT d.pret_id, d.statut_pret, fn_jours_retard(d.pret_id, fn_date_reference()) AS jours_retard
FROM demandes_pret d WHERE d.pret_id IN (9, 61, 67) ORDER BY d.pret_id;
SELECT pret_id, ancien_statut, nouveau_statut, jours_retard, enregistre_par
FROM journal_statuts WHERE enregistre_le = now() ORDER BY journal_id;    -- now() : heure de début de la transaction
ROLLBACK;
\end{lstlisting}

\begin{resultat}
\begin{center}\small
\begin{tabular}{@{}cllc@{}}
\toprule
Prêt & Paiement & Statut après & Retard après \\
\midrule
67 & 916,55~€ : l'échéance du 27/09 est soldée & \textbf{sain} (était en retard) & 0 \\
61 & 300~€ sur l'échéance de 704,41~€ du 15/08 & \code{en_retard} (inchangé) & 46 \\
9 & 9\,422,92~€ : les 13 échéances impayées & \code{en_defaut} (inchangé) & 0 \\
\bottomrule
\end{tabular}
\end{center}
Le journal contient une seule nouvelle ligne : prêt 67, \code{en_retard} $\rightarrow$ \code{sain}.
\end{resultat}

\textbf{Lecture.} Le prêt 61 reste en retard : un paiement partiel ne solde pas l'échéance, dont la date
ne change pas. Le prêt 9 n'a plus aucun retard, mais reste en défaut : c'est la période probatoire,
voulue. Et le filtre \code{enregistre_le = now()} fonctionne parce que \code{now()} renvoie l'heure de
\emph{début} de la transaction, la même pour toutes les instructions.
\end{correction}

% =============================================================================
\section{Exercices à faire seul}

\begin{exercice}{3.4}{Valider l'automatisation contre la réalité, puis faire passer le temps}
\begin{enumerate}[label=\alph*)]
\item Charge le fichier \code{donnees/verite_defauts.csv} (le défaut réel de chaque prêt, issu du jeu de
données) dans une table temporaire. Les prêts que ta logique a mis \og en défaut \fg{} sont-ils
exactement ceux qui ont fait défaut dans la réalité ? Compte les désaccords.
\item Le prêt d'Awa (première échéance le 25/10/2026) : que se passe-t-il au 31/01/2027 si elle ne paie
rien ? Avance la date de référence dans une transaction, regarde son statut avant et après le
traitement de nuit, puis le journal. Annule tout à la fin.
\end{enumerate}
\end{exercice}

\begin{indice}
\begin{itemize}
\item Dans psql : \code{\copy table FROM 'fichier.csv' WITH (FORMAT csv, HEADER true)}. Dans DBeaver :
clic droit sur la table, \emph{Importer des données}.
\item b) Un \code{UPDATE parametres_risque} suffit à « faire passer le temps » : c'est l'intérêt d'avoir
une date de référence paramétrée.
\end{itemize}
\end{indice}

\begin{exercice}{3.5}{La clôture mensuelle des provisions, rejouable sans doublon}
Écris la procédure \code{cloture_mensuelle(date_arrete)}. Elle enregistre dans la table
\code{provisions_mensuelles} la provision de chaque prêt non soldé à cette date (retard, stage, EAD, PD utilisée, LGD, ECL). Elle
refuse une date postérieure à la date de référence et peut être relancée sans créer de doublon.
Produis les 12 arrêtés d'octobre 2025 à septembre 2026, puis l'évolution mensuelle de l'EAD, de l'ECL
par stage, de la \textbf{dotation nette} (variation de l'ECL d'un mois sur l'autre) et du taux de
couverture. Vérifie l'idempotence.
\end{exercice}

\begin{indice}
\begin{itemize}
\item Réutilise \code{fn_situation_prets(date)} : elle sait déjà calculer le passé.
\item \code{CROSS JOIN LATERAL (SELECT ...)} permet de calculer le stage, puis la PD qui en dépend, dans
la même requête.
\item Un \code{CALL} ne s'utilise pas dans un \code{SELECT} : pour boucler sur 12 dates, écris un bloc
anonyme \code{DO $$ ... $$}.
\item La dotation nette : \code{ecl - LAG(ecl) OVER (ORDER BY date_arrete)}.
\end{itemize}
\end{indice}

\begin{exercice}{3.6}{Bonus --- Sécurité : chacun n'a que les droits dont il a besoin}
Crée deux rôles : \code{analyste_risque} (lecture seule sur toutes les tables) et
\code{application_paiement} (aucun droit sur les tables, seulement l'exécution de
\code{enregistrer_paiement}). Rends la procédure \code{SECURITY DEFINER} de façon sûre. Vérifie, dans
une transaction annulée, que l'application peut encaisser par la procédure mais pas écrire directement
dans l'échéancier.
\end{exercice}

\begin{indice}
\begin{itemize}
\item \code{CREATE ROLE ... NOLOGIN} si le rôle n'existe pas (un bloc \code{DO} avec
\code{IF NOT EXISTS (SELECT FROM pg_roles ...)}).
\item \code{REVOKE ALL ON PROCEDURE ... FROM PUBLIC} : par défaut, tout le monde peut exécuter une
procédure.
\item Pour tester : \code{SET LOCAL ROLE application_paiement}, et l'erreur attendue s'appelle
\code{insufficient_privilege}.
\end{itemize}
\end{indice}

% =============================================================================
\section{Synthèse pour les entretiens}

\begin{resume}
\begin{enumerate}
\item On place dans la base ce qui ne doit \textbf{jamais être contourné} : intégrité, traçabilité,
règles réglementaires simples.
\item \textbf{Fonctions} : SQL pour une formule, PL/pgSQL pour une logique ; \code{INTO STRICT} et
\code{RAISE EXCEPTION} pour des erreurs explicites ; volatilité \code{IMMUTABLE}, \code{STABLE},
\code{VOLATILE}.
\item \textbf{Triggers} : \code{BEFORE} pour valider ou compléter, \code{AFTER} pour réagir ; par ligne ou
par instruction (tables de transition) ; \code{WHEN} pour filtrer ; attention aux cascades.
\item \textbf{Procédures} : \code{CALL}, opérations métier, verrou \code{FOR UPDATE} contre les accès
simultanés.
\item \textbf{Le temps n'est pas un événement} : un traitement planifié complète les triggers, avec la
même règle.
\item \textbf{Idempotence} : \code{ON CONFLICT DO UPDATE}. \textbf{Alertes} : \code{pg_notify}, délivrées
au commit. \textbf{Sécurité} : moindre privilège, \code{SECURITY DEFINER} avec \code{search_path} figé.
\end{enumerate}
\end{resume}

\begin{entretien}
\begin{itemize}
\item « La détection du défaut est automatique et validée : sur 32\,409 prêts réels, ma logique trouve
7\,088 défauts, exactement les 7\,088 du jeu de données, sans aucun désaccord. »
\item « J'ai séparé les responsabilités : un trigger \code{BEFORE} protège le contrat, un trigger
\code{AFTER} par instruction recalcule les statuts en un seul passage grâce aux tables de transition,
un troisième journalise et alerte. La règle du statut est une fonction, partagée avec le traitement de
nuit. »
\item « Un trigger ne voit pas le temps passer : sans paiement, un retard franchit 90 jours sans
événement. C'est pourquoi un traitement planifié complète les triggers. »
\item « La clôture mensuelle est idempotente grâce à \code{ON CONFLICT} : on peut la relancer sans
doublon. Et l'application d'encaissement n'a aucun droit sur les tables : elle passe par une procédure
\code{SECURITY DEFINER} au \code{search_path} figé. »
\end{itemize}
\end{entretien}

\subsection*{Questions fréquentes en entretien}

\begin{description}[leftmargin=0pt, itemsep=0.5em]
\item[Fonction ou procédure ?] Une fonction calcule et renvoie une valeur, utilisable dans un
\code{SELECT} ; une procédure exécute une opération, s'appelle par \code{CALL} et peut gérer les
transactions.
\item[BEFORE ou AFTER ?] \code{BEFORE} pour valider ou modifier la ligne avant écriture ; \code{AFTER}
pour réagir au résultat (journal, recalcul d'une autre table).
\item[Quels risques avec les triggers ?] Logique invisible, cascades, récursion, coût par ligne, et une
erreur qui annule toute l'instruction.
\item[Comment éviter qu'un même paiement soit compté deux fois ?] Verrouiller la ligne du prêt
(\code{FOR UPDATE}) et contraindre les montants (\code{paiement_borne}).
\item[Qu'est-ce que l'idempotence ?] Une opération qu'on peut relancer sans changer le résultat : ici,
une clôture rejouée remplace ses lignes au lieu de les dupliquer.
\end{description}
